\chapter{Noise2Noise : a self-supervised framework}


In this section, we present several denoising methods for a certain signal $x$ corrupted by noise to produce a noisy observation $y$.
Let $(\Omega, \mathcal{F}, \mathbb{P})$ be a probability space.
Let $X : \Omega \to \mathbb{R}^n$ and $Y : \Omega \to \mathbb{R}^n$ be random variables representing respectively the clean signal and the noisy observation.
We consider a corruption process $\mathcal{C}$ that maps the clean signal $X$ to the noisy observation $Y$ according to a certain noise model. This corruption process can be expressed as follows :

\begin{equation}
    Y | X=x \sim \mathcal{C}(x)
    \label{eq:noise-model}
\end{equation}

We will consider the residual additive noise $E$ defined as follows :

\begin{equation}
    E = Y - X
\end{equation}

\section{Self-supervised learning for denoising}

\subsection{Denoising Auto-Encoders}

Whereas traditional auto-encoders are trained to reconstruct their input through an encoder-decoder structure, Denoising auto-encoders (DAEs) are trained to reconstruct clean data from corrupted versions of it.
They have been first introduced by \cite{vincent_extracting_2008} and have since been widely used in various applications such as image denoising, inpainting or other inverse problems.
The DAE objective is thus :

\begin{equation}
    \mathcal{L}_{DAE}(\theta) = \underset{\substack{X \sim p_X \\ Y | X=x \sim \mathcal{C}(x)}}{\mathbb{E}} \left[ \| X - f_{\theta}(Y) \|_2^2 \right]
\end{equation}
with $\mathcal{C}$ being a random corruption process.

\subsection{Score-based denoising}


\begin{definition}{Score function}
    % The score function of a random variable $X$ with $\mathcal{C}^1$ probability density function $p(x)$ is defined as the gradient of the log-density with respect to $x$ :
    \begin{equation}
        s_X(x) = \nabla_x \ln p(x)
    \end{equation}
\end{definition}

This function denotes the direction of maximum increase of the probability density function at point $x$.

\subsubsection{Score matching estimation}

As first introduced by \cite{hyvarinen_estimation_2005}, score matching is a method to estimate parametric models of probability density functions without requiring the computation of the normalizing constant.
Given a parametric model score function $s(\cdot; \theta)$ with parameter $\theta$, and the data score function $s_X(\cdot)$ of observed data sampled from $X$, we want to minimize the distance between the data score function $s_X(\cdot)$ and the model score function $s_X(\cdot; \theta)$ using the following objective :

\begin{equation}
    \mathcal{L}_{SM}(\theta) = \mathbb{E}_X \left[ \| s_X(x) - s(x; \theta) \|_2^2 \right]
\end{equation}

Minimization of this objective does not require knowledge of the normalizing constant of $s(\cdot; \theta)$ as it only involves derivatives of the log-density.
They also show that we don't need to know the data score function $s_X(x)$ to optimize this objective as it can be rewritten as :

\begin{equation}
    \mathcal{L}_{SM}(\theta) = \mathbb{E}_X \left[ \mathop{Tr} \left( \nabla_x s(x; \theta) \right) + \frac{1}{2} \| s_X(x; \theta) \|_2^2 \right] + C
\end{equation}

This is very convenient as it avoids the need to estimate the data score function over the observed data, which is generally intractable.

\subsubsection{Denoising Auto-Encoders and score function estimation}




DAEs and score matching principles have been linked together by \cite{vincent_connection_2011}. They built an objective similar to score matching but defined on clean/noisy image pairs with the corruption process $\mathcal{C}$ distribution being known.

\begin{equation}
    \mathcal{L}_{DSM}(\theta) = \underset{X \sim p_X \\ Y \sim \mathcal{C}(X)}{\mathbb{E}} \left[ \| s_X(Y, \theta) - s_{Y|X}(Y|X) \|_2^2 \right]
\end{equation}

They prove the specific case of Tweedie's formula for additive Gaussian noise (see Theorem \ref{theorem:tweedie_Gaussian}) which will be considered in the following.

\cite{alain_what_2014} further prove that when the Gaussian noise level $\sigma$ tends to zero, the DAE learns the score of the data distribution as follows :

\begin{equation}
    f_{\theta^*}(y) = y + \sigma^2 s_X(y) + \underset{\sigma \to 0}{o}(\sigma^2)
\end{equation}

This can be generalized to a wider range of exponential family distributions thanks to Tweedie's formula \cite{efron_tweedies_2011}. We retain the two following theorems that stand for additive Gaussian noise and Poisson noise respectively.

\begin{theorem}{Tweedie's formula for additive Gaussian noise}
    \label{theorem:tweedie_Gaussian}
    Let $X$ be a random variable and $Y = X + N$ be a noisy observation of $X$ corrupted by an additive Gaussian noise $N$ independent from $X$ with variance $\sigma^2$. The MMSE estimator of $X$ given $Y = y$ is given by :
    \begin{equation}
        \mathbb{E}\left[X | Y = y\right] = y + \sigma^2 \nabla_y \ln p(y)
    \end{equation}
\end{theorem}

% \begin{proof}
%     Bayes' rule gives us:
%     \begin{equation*}
%         p(x|y) = \frac{p(y|x)p(x)}{p(y)}
%     \end{equation*}
%     The likelihood $p(y|x)$ is given by the noise distribution:
%     \begin{equation*}
%         p(y|x) = \frac{1}{\sqrt{2 \pi \sigma^2}} \exp\left(-\frac{||y - x||^2}{2 \sigma^2}\right)
%     \end{equation*}
%     Thus we have:
%     \begin{equation*}
%         p(x|y) = \frac{1}{Z(y)} \exp\left(-\frac{||y - x||^2}{2 \sigma^2}\right) p(x)
%     \end{equation*}
%     with $Z(y)$ acting as a normalizing constant:
%     \begin{equation*}
%         Z(y) = \int p(x) \exp\left(-\frac{||y - x||^2}{2 \sigma^2}\right)dx
%     \end{equation*}
%     Then we compute the derivative of $p(y)$:
%     \begin{align*}
%         \nabla_y p(y) &= \frac{d}{dy} \int p(y|x)p(x) dx \\
%         &= \int \frac{\partial}{\partial y} p(x) \exp\left(-\frac{||y - x||^2}{2 \sigma^2}\right) dx \\
%         &= - \frac{1}{\sigma^2} \int (y - x) p(x) \phi_{\sigma}(y - x) dx \\
%     \end{align*}
%     where we denote $\phi_{\sigma}$ the Gaussian kernel of standard deviation $\sigma$.

%     \begin{align*}
%         \nabla_y \ln p(y) = \frac{p'(y)}{p(y)} &= - \frac{1}{\sigma^2} \left( \int (y - x) \frac{p(x) \phi_{\sigma}(y - x)}{p(y)} dx \right) \\
%         &= - \frac{1}{\sigma^2} \left( y - \int x p(x|y) dx \right) \\
%         &= - \frac{1}{\sigma^2} (y - \mathbb{E}[x|y])
%     \end{align*}
% \end{proof}

\begin{theorem}{Tweedie's formula for Poisson noise}
    Let $X$ be a random variable and $Y$ be a noisy observation of $X$ corrupted by Poisson noise. The MMSE estimator of $X$ given $Y = y$ is given by :
    \begin{equation}
        \mathbb{E}\left[X | Y = y\right] = (y + 1)\exp \left( \nabla_y \ln p(y)\right)
    \end{equation}
\end{theorem}

% \begin{proof}
%     The assumption of Poisson noise means that the likelihood $p(y|x)$ is given by:
%     \begin{equation*}
%         p(y|x) = \frac{(x)^{y} e^{-x}}{y!}
%     \end{equation*}
%     Then we compute the gradient of $p(y)$:
%     \begin{align*}
%         \exp(\ln \nabla_y p(y)) &= \frac{p(y+1)}{p(y)} \\
%         &= \frac{1}{p(y)} \int p(x) \frac{x}{y+1}\frac{x^y}{y!}e^{-x} dx \\
%         &= \frac{1}{y+1} \int x \frac{p(x) p(y|x)}{p(y)} dx \\
%         &= \frac{1}{y+1} \mathbb{E}[x|y]
%     \end{align*}
%     Thus we have:
%     \begin{equation*}
%         \mathbb{E}[x|y] = (y+1) \exp(\ln \nabla_y p(y))
%     \end{equation*}
% \end{proof}


Given the previous proposition and theorems, one can build a link between denoising auto-encoders and score function estimation as follows :

\begin{equation}
    s_Y(y) = \frac{D_{\sigma}(y) - y}{\sigma^2} \quad \text{(Gaussian noise)}
\end{equation}
\begin{equation}
    s_Y(y) = \ln \left( \frac{D_{Poisson}(y)}{y + 1} \right) \quad \text{(Poisson noise)}
\end{equation}

We can highlight an intriguing application of \cite{kim_noise2score_nodate} who leverage this link to estimate the score function using a DAE trained under an additive Gaussian noise assumption, even when the actual noise model is not Gaussian. They then use the estimated score function to denoise images corrupted by other types of noise using the appropriate variant of Tweedie's formula.

\subsection{Stein's Unbiased Risk Estimator}

\section{The Noise2Noise framework}

A seminal work in the denoising field is the \textbf{Noise2Noise} framework \cite{lehtinen_noise2noise_2018}.
Noise2Noise is a denoising method at the crossroad between self-supervised methods — which learn to denoise only from noisy samples — and supervised methods — which require clean targets to learn the mapping between noisy and clean image.
The key idea of Noise2Noise is that we learn to denoise images from pairs of noisy images instead of noisy-clean image pairs or single noisy images.
These pairs of noisy images can be obtained in various ways depending on the application. Each pair must represent the same underlying clean image but with different realizations of the noise process.
Hence, we use the information contained in the noisy data itself to learn how to denoise it, without ever observing clean images.

Noise2noise can be extended to the case of linear inverse problems in a more generic setting. We will present this extension in the last part of this section. As a consequence, we will use the formalism of
\eqref{eq:inverse-problem} to present the Noise2Noise framework and its theoretical guarantees.

Let's first consider the most basic case of pure denoising, i.e. when $A = I$ in \eqref{eq:inverse-problem}.
In order to avoid confusion, we will not write $Y = AX + B$ but rather $Y = X + B$ in the following.

Let $Y^{(1)}$ and $Y^{(2)}$ be two noisy observations of the same underlying clean image $X$ obtained through the same corruption process $\mathcal{C}$ but with different realizations of the noise.
Thus, the noisy pairs are generated according to the following noise model :

\begin{equation}
    Y^{(j)} | X \sim \mathcal{C}(X), \quad j = 1, 2
    \label{eq:additive_noise}
\end{equation}

Let $\{(Y_i^{(1)}, Y_i^{(2)})\}_{i=1}^N$ be a finite sequence of independant and statistically distributed random variables.
We denote as $y_i^{(j)}$ the observed noisy images for $j = 1, 2$ and $i = 1, \dots, N$.

We train a neural network $f_\theta$ with parameters $\theta$ to map the noisy image $y_i^{(1)}$ to the other noisy image $y_i^{(2)}$ both representing a different corruption of the same underlying clean image $x_i$ and vice versa by minimizing the following loss function :

\begin{equation}
    \mathcal{L}_{\text{N2N}}(\theta) = \frac{1}{2N} \sum_{i=1}^N \left[ \rho(f_\theta(y_i^{(1)}), y_i^{(2)}) + \rho(f_\theta(y_i^{(2)}), y_i^{(1)}) \right]
    \label{eq:n2n_loss}
\end{equation}
with $\rho$ a loss function such as the Mean Squared Error (MSE), the Mean Absolute Error (MAE) or any other appropriate loss.

The supervised counterparts of \eqref{eq:n2n_loss} would be :

\begin{equation}
    \mathcal{L}_{\text{clean}}(\theta) = \frac{1}{N} \sum_{i=1}^N \rho(f_\theta(y_i), x_i) 
    \label{eq:supervised_loss}
\end{equation}
with $y_i$ being a noisy observation of the clean image $x_i$.

One can have the intuition that the loss \eqref{eq:n2n_loss} is not equivalent to the supervised loss \eqref{eq:supervised_loss} as we do not have access to the clean image $x_i$.
Intuitively, a Noise2Noise-trained model will learn to push a sample toward the mean of the distribution of the noisy samples as illustrated in Figure \ref{fig:noise2noise}.

\begin{figure}
    \centering
        \begin{tikzpicture}
            % Network prediction node
            \node [draw, circle, fill=blue!40, thick] (pred) {\textbf{$\hat{y}$}};
            \node [above=0.15cm] at (pred.north) {Network prediction};

            % Clean target node (not directly used, but for reference)
            \node [shade, inner color=red!40, outer color=white, draw=none] (clean_target_halo) at ([xshift=-0.25cm, yshift=-1.5cm]pred.south west) [circle, minimum size=1.5cm] {};
            \node [draw, circle, fill=red!40, thick] (clean_target) at (clean_target_halo)  {\textbf{$y$}};
            \node [below=0.1cm] at (clean_target.south) {Unobserved target};

            % Multiple noisy targets
            \foreach \i/\angle/\distance in {1/210/2.3, 2/125/1.5, 3/290/1.8, 4/0/2, 5/130/3} {
                \node[draw, circle, fill=orange!30, thick] (noisy\i) at ($(clean_target) + (\angle:\distance)$) {\textbf{$y_{\i}$}};
                \draw[->, thick, opacity=0.4] (pred) -- (noisy\i);
            }
        \end{tikzpicture}
    \caption{Illustration of the Noise2Noise framework where a neural network is trained to map a noisy input image to multiple noisy target images representing different realizations of the same underlying clean image.}
    \label{fig:noise2noise}
\end{figure}

\subsection{Specific settings for Noise2Noise}

\paragraph{MSE loss case}

\begin{proposition}
    \label{prop:noise2noise_l2}
    Minimizing \eqref{eq:n2n_loss} with $\rho(x, y) = \| x - y \|_2^2$ gives the same estimator as minimizing \eqref{eq:supervised_loss} as $N \to \infty$ under the following conditions :
    \begin{itemize}
        \item $Y^{(1)}$ and $Y^{(2)}$ are independent given the clean image $X$ (i.e. $Y^{(1)} | X \independent Y^{(2)} | X$).
        \item Noise is centered, i.e. $\mathbb{E}[Y^{(j)} | X^{(j)} = x] = x$ for $j = 1, 2$.
    \end{itemize}
\end{proposition}

\begin{proof}
    We expand the Noise2Noise empirical risk as follows :
    \begin{align*}
        R_{N,1}^{n2n}(\theta) &= \frac{1}{N} \sum_{i=1}^N \left[ \| f_\theta(Y_i^{(1)}) - (X^{(i)} + B^{(2)}) \|_2^2 \right] \\
        &= \frac{1}{N} \sum_{i=1}^N \left[ \|( f_\theta(Y_i^{(1)}) - X^{(i)}) - B^{(2)} \|_2^2 \right] \\
        &= R_{N,1}^{clean}(\theta) + \frac{1}{N} \sum_{i=1}^N \left[ \| B^{(2)} \|_2^2 \right] - \frac{2}{N} \sum_{i=1}^N \left[ (f_\theta(Y_i^{(1)}) - X^{(i)})^T B^{(2)} \right]
    \end{align*}
    However, $\frac{2}{N} \sum_{i=1}^N \left[ (f_\theta(Y_i^{(1)}) - X^{(i)})^T B^{(2)} \right] \xrightarrow[N \to \infty]{a.s.} 2 \mathbb{E} \left[ (f_\theta(Y^{(1)}) - X)^T B^{(2)} \right]$ by the law of large numbers. This expectation vanishes under the independence and zero-mean assumptions on the noise.
    Thus, we have :
    \begin{equation*}
        R_{N,1}^{n2n}(\theta) = R_{N,1}^{clean}(\theta) + \frac{1}{N} \sum_{i=1}^N \left[ \| B^{(2)} \|_2^2 \right] + o(1)
    \end{equation*}
    We show a similar result for the second term of the Noise2Noise empirical risk by swapping $Y^{(1)}$ and $Y^{(2)}$ and we get the following result :
    \begin{equation*}
        R_{N}^{n2n}(\theta) = 2 R_{N}^{clean}(\theta) + \frac{1}{N} \sum_{i=1}^N \left[ \| B_i^{(1)} \|_2^2 + \| B_i^{(2)} \|_2^2 \right] + o(1)
    \end{equation*}
    where we are left with the standard clean denoising empirical risk with its underlying bias and variance decomposition, and an irreductible noise variance term which has no influence as it is independent of $\theta$.
\end{proof}

\subsubsection{MSE Noise2Noise loss for aggregated outputs}

The use of the MSE loss function enables the use of an even more suited training objective for the case of aggregated outputs.
The two independant noisy observations may be obtained by splitting the data into two
different subsets, hence the noise in each subset is higher than the noise in the original data.
This is the case for example in the medical field.
\cite{wu19} propose a consensus loss to alleviate this issue and improve the results of the aggregated outputs 
\begin{equation}
    \hat{y_i} = \frac{1}{2}(f_\theta(y_i^{(1)}) + f_\theta(y_i^{(2)}))
\end{equation}

We start by stating that the supervised loss of the aggregated output can be written as follows :

\begin{align*}
    \mathcal{L}_{\text{clean}}(\theta) &= \frac{1}{N} \sum_{i=1}^N \| \frac{1}{2}(f_\theta(y_i^{(1)}) + f_\theta(y_i^{(2)})) -  x_i \|_2^2 \\
    &= \frac{1}{N} \sum_{i=1}^N \frac{1}{2} \| f_{\theta}(y_i^{(1)}) - x_i \|_2^2 + \frac{1}{2} \| f_{\theta}(y_i^{(2)}) - x_i \|_2^2 - \frac{1}{4} \| f_{\theta}(y_i^{(1)}) - f_{\theta}(y_i^{(2)}) \|_2^2
\end{align*}

Theorem \ref{prop:noise2noise_l2} shows that the the unobserved clean image $x_i$ can be replaced by one of its noisy version $y_i^{(j)}$ in the previous expression without changing the optimal solution of the problem.
Thus, we can rewrite the supervised loss of the aggregated output as follows :

\begin{equation}
    \mathcal{L}_{\text{consensus}}(\theta) = \frac{1}{N} \sum_{i=1}^N \frac{1}{2} \| f_{\theta}(y_i^{(1)}) - y_i^{(2)} \|_2^2 + \frac{1}{2} \| f_{\theta}(y_i^{(2)}) - y_i^{(1)} \|_2^2 - \frac{1}{4} \| f_{\theta}(y_i^{(1)}) - f_{\theta}(y_i^{(2)}) \|_2^2
    \label{eq:consensus_loss}
\end{equation}

\paragraph{Poisson negative log-likelihood case}
\label{sec:noise2noise_poisson}

As we have a strong interest in PET imaging, we will also focus on the case of Poisson noise which is the best model for PET data and other photon-limited imaging modalities.

\begin{definition}
    For two random variables $X: \Omega \to \mathcal{X}$ and $Y: \Omega \to \mathcal{Y}$, we say that $Y \,|\, X \sim \text{Poisson}(X)$ if :
    \begin{equation*}
        \forall i = 1 \dots d, \quad \mathbb{P}(Y[i] = k \,|\, X[i] = x[i]) = \frac{(x[i])^k e^{-x[i]}}{k!}, \quad k \in \mathbb{N}
    \end{equation*}
\end{definition}

\begin{proposition}
    Under the assumption $Y \,|\, X \sim \text{Poisson}(X)$, we have :
    \begin{equation*}
        \mathbb{E}[Y \,|\, X] = X
    \end{equation*}
\end{proposition}

\begin{proof}
    By definition of the Poisson distribution, we have :
    \begin{align*}
        \mathbb{E}[Y[i] \,|\, X[i] = x[i]] &= \sum_{k=0}^{\infty} k \cdot \mathbb{P}(Y[i] = k \,|\, X[i] = x[i]) \\
        &= \sum_{k=0}^{\infty} k \cdot \frac{(x[i])^k e^{-x[i]}}{k!} \\
        &= x[i] \sum_{k=1}^{\infty} \frac{(x[i])^{k-1} e^{-x[i]}}{(k-1)!} \\
        &= x[i] \cdot 1 = x[i]
    \end{align*}
    Thus, we have $\mathbb{E}[Y \,|\, X] = X$.
\end{proof}

Now, we consider the case of Poisson noise (i.e. $Y^{(j)} \,|\, X^{(j)} \sim \text{Poisson}(X^{(j)}), \quad j = 1, 2$).
Given a clean image $x$ and a noisy observation $y$ corrupted by Poisson noise, the likelihood of observing $y$ given $x$ is given by :
\begin{equation*}
    p(y|x) = \prod_{i} \frac{(x[i])^{y[i]} e^{-x[i]}}{y[i]!}
\end{equation*}
Taking the negative logarithm of this likelihood gives us the Poisson negative log-likelihood loss as defined in the following definition.

\begin{definition}
    The Poisson negative log-likelihood loss between two images $y$ and $x$ is defined as follows :
    \begin{equation*}
        \mathcal{L}_{Poisson}(y, x) = \sum_{i} \left[ y[i] - x[i] \log(y[i]) + \log(x[i]!) \right]
    \end{equation*}
\end{definition}

\begin{proposition}
    \label{prop:noise2noise_poisson}
    If $Y | X=x \sim \text{Poisson}(x)$, then minimizing \eqref{eq:n2n_loss} with $\rho = \mathcal{L}_{Poisson}$ gives the same estimator as minimizing \eqref{eq:supervised_loss} as $N \to \infty$ under the following conditions :
    \begin{itemize}
        \item $Y^{(1)}$ and $Y^{(2)}$ are independent given the clean image $X$ (i.e. $Y^{(1)} | X \independent Y^{(2)} | X$)..
    \end{itemize}
\end{proposition}

\begin{proof}
    The poisson negative log-likelihood between $Y_1$ and $Y_2$ can be expressed as follows :
    \begin{equation*}
        \mathcal{L}_{Poisson}(Y_1, Y_2) = Y_1 - Y_2 \log(Y_1) + \log(Y_2!)
    \end{equation*}
    We expand the Noise2Noise Poisson empirical risk as follows :
    \begin{align*}
        R_{N,1}^{n2n}(\theta) &= \frac{1}{N} \sum_{i=1}^N \left[ \mathcal{L}_{Poisson}(f_\theta(Y_i^{(1)}), Y_i^{(2)}) \right] \\
        &= \frac{1}{N} \sum_{i=1}^N \left[ f_\theta(Y_i^{(1)}) - Y_i^{(2)} \log(f_\theta(Y_i^{(1)})) + \log(Y_i^{(2)}!) \right]
    \end{align*}
    Yet, $\frac{1}{N} \sum_{i=1}^N Y_i^{(2)} \log(f_\theta(Y_i^{(1)})) \xrightarrow[N \to \infty]{a.s.} \mathbb{E} \left[ Y^{(2)} \log(f_\theta(Y^{(1)})) \right]$ due to the law of large numbers. Under the independence assumption on the noise, this expectation can be rewritten as 
    \begin{align*}
        \mathbb{E} \left[ Y^{(2)} | X \right] \mathbb{E} \left[ \log(f_\theta(Y^{(1)})) \right] = X \mathbb{E} \left[ \log(f_\theta(Y^{(1)})) \right]
    \end{align*}
    Then, the empirical risk satisfies the following almost sure convergence :
    \begin{align*}
        R_{N,1}^{n2n}(\theta) \xrightarrow[N \to \infty]{a.s.} \mathbb{E} \left[ f_\theta(Y^{(1)}) - X \log(f_\theta(Y^{(1)})) + \log(Y^{(2)}!) \right]
    \end{align*}
    We can do the same for the second term of the loss by swapping $Y^{(1)}$ and $Y^{(2)}$ and we get the following result :
    \begin{equation*}
        R_{N}^{n2n}(\theta) \xrightarrow[N \to \infty]{a.s.} 2 \mathbb{E} \left[ f_\theta(Y) - X \log(f_\theta(Y)) + \log(Y!)\right]
    \end{equation*}
    where we are left with the standard clean poisson negative-log likelihood risk plus a constant term independent of $\theta$.
\end{proof}

\subsection{General Noise2Noise theory}

In proofs of Proposition \ref{prop:noise2noise_l2} and \ref{prop:noise2noise_poisson}, we can see that the empirical risks of the Noise2Noise loss functions are asymptotically the same
up to an additional term independent of the parameters $\theta$ to be optimized, hence leading to the same optimal solution.
In this section, we try to give more general theoretical insights to understand the Noise2Noise framework and its guarantees.
The choice of the loss function $\rho$ is crucial to ensure those guarantees and we will see that the noise model also plays a key role in the effectiveness of the Noise2Noise framework.

\subsubsection{Noise-induced loss bias}

\begin{definition}
    We define the noise-induced loss bias as between two vectors $z$ and $x$ as the difference between the expected value of the loss function $\rho$ evaluated at $z$ and a noisy observation $Y$ and the actual distance between $z$ and $x$, i.e. :
    \begin{equation}
        \Delta_\rho(z, x) = \mathbb{E}[\rho(z, Y) | X=x] - \rho(z, x)
    \end{equation}
\end{definition}

We denote as $R^{\text{N2N}}(\theta)$ (resp. $R^{\text{clean}}(\theta)$) the expected value of the Noise2Noise loss (resp. the supervised loss) over the distribution of the noisy observations.

\begin{equation}
    R^{\text{N2N}}(\theta) = \mathbb{E}_{X, Y^{(1)}, Y^{(2)}}[\rho(f_\theta(Y^{(1)}), Y^{(2)})]
\end{equation}

\begin{equation}
    R^{\text{clean}}(\theta) = \mathbb{E}_{X, Y}[\rho(f_\theta(Y), X)]
\end{equation}

\begin{proposition}
    The Noise2Noise and the supervised expected risks are tied by the following relation :
    \begin{equation}
        R^{\text{N2N}}(\theta) = R^{\text{clean}}(\theta) + \mathbb{E}[\Delta_\rho(f_\theta(Y), X)]
    \end{equation}
\end{proposition}

\begin{proof}
    We have :
    \begin{align*}
        R^{\text{N2N}}(\theta) &= \mathbb{E}_{X, Y^{(1)}, Y^{(2)}}[\rho(f_\theta(Y^{(1)}), Y^{(2)})] \\
        &= \mathbb{E}_{X, Y^{(1)}}[\mathbb{E}_{Y^{(2)}}[\rho(f_\theta(Y^{(1)}), Y^{(2)}) | X, Y^{(1)}]] \\
        &= \mathbb{E}_{X, Y^{(1)}}[\mathbb{E}_{Y^{(2)}}[\rho(f_\theta(Y^{(1)}), Y^{(2)}) | X]]
    \end{align*}
    Let $z = f_\theta(Y^{(1)})$, we have $\mathbb{E}_{Y^{(2)}}[\rho(z, Y^{(2)}) | X] = \rho(z, X) + \Delta_\rho(z, X)$ by definition of the noise-induced loss bias. Thus, we have :
    \begin{align*}
        R^{\text{N2N}}(\theta) &= \mathbb{E}_{X, Y^{(1)}}[\rho(f_\theta(Y^{(1)}), X) + \Delta_\rho(f_\theta(Y^{(1)}), X)] \\
        &= R^{\text{clean}}(\theta) + \mathbb{E}[\Delta_\rho(f_\theta(Y), X)]
    \end{align*}
\end{proof}

The $\mathbb{E}[\Delta_\rho(f_\theta(Y), X)]$ term is very informative regarding the effectiveness of the Noise2Noise framework for a given noise model and loss function.
If this term is independent of the parameters $\theta$, then minimizing the Noise2Noise loss is equivalent to minimizing the supervised loss and we can learn to denoise without ever observing clean images.

\begin{corollary}
    \label{coro:delta}
    The minimization of $R^{\text{N2N}}(\theta)$ yields the same parameters as the minimization of $R^{\text{clean}}(\theta)$ if and only if $\Delta_\rho(z, x)$ is independent of $z$.
\end{corollary}

\subsubsection{Noise2Noise for centered noise}

\begin{definition}
    We say that a noise model is centered if the expected value of the noisy observation $Y$ given the clean image $X=x$ is equal to $x$, i.e. $\mathbb{E}[Y | X=x] = x$.
\end{definition}

\begin{proposition}
    \label{prop:noise2noise_centered_mse}
    If the noise is centered, i.e. $\mathbb{E}[Y | X=x] = x$, then minimizing the Noise2Noise loss with the MSE loss function is equivalent to minimizing the supervised loss as $N \to \infty$.
\end{proposition}

\begin{proof}
If the MSE loss function is used, the noise-induced loss bias is calculated as follows :
\begin{align*}
    \Delta_\rho(z, x) &= \mathbb{E}[\| z - Y \|_2^2 | X=x] - \| z - x \|_2^2 \\
    &= \mathbb{E}[\| z - x + (x - Y) \|_2^2 | X=x] - \| z - x \|_2^2 \\
    &= \mathbb{E}[\| z - x \|_2^2 + \| x - Y \|_2^2 + 2 \langle z - x, x - Y \rangle | X=x] \\
    &\quad - \| z - x \|_2^2 \\
    &= \mathbb{E}[\| x - Y \|_2^2 | X=x] + 2 \langle z - x, \mathbb{E}[x - Y | X=x] \rangle
\end{align*}
According to corollary \ref{coro:delta}, we need this quantity to be independent of $z$. Thus, we need the term $\langle z - x, \mathbb{E}[x - Y | X=x] \rangle$ to vanish for all $z$, which is equivalent to having $\mathbb{E}[Y | X=x] = x$.
\end{proof}

Note that we recover the result of Proposition \ref{prop:noise2noise_l2} in the case of additive centered noise and the MSE loss function.
We have proven that the Noise2Noise framework, combined with the use of the MSE loss function, is valid for any noise model as long as the noise is centered around its clean counterpart,
in addition to the independence assumption between the two noisy observations given the clean image.
For example, if $Y|X=x$ follows a log-normal distribution, then $\mathbb{E}[Y|X=x] = x \exp(\sigma^2/2) \neq x$ and the Noise2Noise framework with the MSE loss function is not valid in this case.

\subsubsection{Toward the exponential family of distributions}
As we already saw in the case of the MSE loss under gaussian noise and the Poisson negative log-likelihood loss under Poisson noise, the additional term $\mathbb{E}[\Delta_\rho(f_\theta(Y), X)]$ is respectively equal to the variance of the noise and a constant term independent of $\theta$.
We can show that this property can be extended under some asumptions to the exponential family of distributions and for the corresponding negative log-likelihood loss function.
\begin{definition}
    A distribution belongs to the exponential family if its probability density function can be written as :
    \begin{equation}
        p(y|x) = h(y) \exp(\langle T(y), \eta(x) \rangle - \phi(x))
    \end{equation}
    where $T(y)$ is the sufficient statistic, $\eta(x)$ is the natural parameter, $\phi(x)$ is the log-partition function and $h(y)$ is the base measure.
\end{definition}
Using the exponential family formalism, we can write the negative log-likelihood loss function as follows :
\begin{equation}
    \rho(z, y) = -\log p(y|z) = -\log h(y) - \langle T(y), \eta(z) \rangle + \phi(z)
\end{equation}
Then, computing the noise-induced loss bias gives us :
\begin{align*}
    \Delta_\rho(z, x) &= \mathbb{E}[\rho(z, Y) | X=x] - \rho(z, x) \\
    &= \mathbb{E}[-\log h(Y) - \langle T(Y), \eta(z) \rangle + \phi(z) | X=x] - (-\log h(x) - \langle T(x), \eta(z) \rangle + \phi(z)) \\
    &= \mathbb{E}[-\log h(Y) | X=x] + \log h(x) - \langle \mathbb{E}[T(Y) | X=x] - T(x), \eta(z) \rangle + \phi(z) - \phi(z) \\
    &= \mathbb{E}[-\log h(Y) | X=x] + \log h(x) - \langle \mathbb{E}[T(Y) | X=x] - T(x), \eta(z) \rangle
\end{align*}
For this quantity to be independent of $z$, we need the term $\langle \mathbb{E}[T(Y) | X=x] - T(x), \eta(z) \rangle$ to vanish, which is equivalent to having $\mathbb{E}[T(Y) | X=x] = T(x)$, hence Proposition \ref{prop:ef_n2n}.
\begin{proposition}
    \label{prop:ef_n2n}
    For a noise model belonging to the exponential family, if the sufficient statistic $T(Y)$ is an unbiased estimator of $T(X)$, then minimizing the Noise2Noise loss with the corresponding negative log-likelihood loss function is equivalent to minimizing the supervised loss as $N \to \infty$.
\end{proposition}
Note that this property of the sufficient statistics is satisfied for Gaussian and Poisson distributions as we have already seen in Proposition \ref{prop:noise2noise_l2} and \ref{prop:noise2noise_poisson}
but also for other distributions like the Bernoulli distribution, the Exponential distribution, or the Gamma distribution.
On the other hand, the log-normal distribution belongs to the exponential family but its sufficient statistic is $T(Y) = (\log Y, (\log Y)^2)$.
$\mathbb{E}[\log Y | X=x] = \log x - \frac{\sigma^2}{2}$ and $\mathbb{E}[(\log Y)^2 | X=x] = (\log x)^2 + \sigma^2 - \sigma^2 \log x + \frac{\sigma^4}{4}$, hence the sufficient statistic is not an unbiased estimator of $T(X) = (\log X, (\log X)^2)$.

\subsubsection{Empirical Noise2Noise}

The two previous paragraphs show that neither the MSE loss function nor the log-normal negative log-likelihood loss function are appropriate for the Noise2Noise framework under log-normal noise.
For log-normal distributions or any other noise model that doesn't satisfy either the centered noise assumption or the sufficient statistic unbiasedness assumption,
we can empirically observe that the Noise2Noise framework will indeed tend to denoise the images but will not converge to the supervised solution as the number of training samples increases.
In such cases, we lack interpretability as the resulting estimator is no longer the optimal solution of a well-defined supervised loss function or even a MMSE estimator in some cases.

\subsubsection{Particular case of the MAE loss}

The Minimum Absolute Error (MAE) loss function is defined as $\rho(z, y) = \| z - y \|_1$ is a special case or the Noise2Noise framework.
Its noise-induced loss bias is given by $\Delta_\rho(z, x) = \mathbb{E}[\| z - Y \|_1 | X=x] - \| z - x \|_1$.
If $g_x(z) = \mathbb{E}[\| z - Y \|_1 | X=x]$, then $g_x'(z) = 2 \mathbb{P}(Y \leq z | X=x) - 1 = 2 F_{Y|X=x}(z) - 1$ where $F_{Y|X=x}$ is the cumulative distribution function of $Y$ given $X=x$.
In order to apply Corollary \ref{coro:delta}, we need $g_x'(z) = 0$ for all $z$, which is equivalent to having $F_{Y|X=x}(z) = 1/2$ for all $z$, which is impossible.
Thus, the MAE loss function is not appropriate for the Noise2Noise framework ; whereas the supervised loss with the MAE loss function will converge to the conditional median of $Y$ given $X=x$, the Noise2Noise loss will not converge to any well-defined estimator.

\subsection{Extension to Bregman Geometry}

\subsubsection{Bregman Geometry}

In this section, we study the potential Noise2Noise guarantees in the context of Bregman geometry, which is a generalization of the Euclidean geometry, the MSE loss function or any other loss function that can be expressed as a Bregman divergence.
\begin{definition}
Let $\phi: \mathbb{R}^d \to \mathbb{R}$ be a strictly convex and differentiable function. The Bregman divergence associated with $\phi$ is defined as :
\begin{equation}
    \label{bregman_divergence}
    D_\phi(z, y) = \phi(z) - \phi(y) - \langle \nabla \phi(y), z - y \rangle
\end{equation}
\end{definition}
\noindent The Bregman divergence is not a distance in the mathematical sense as it is not symmetric and does not satisfy the triangle inequality, but it is a measure of dissimilarity between two points in the space defined by $\phi$.
It is though positive and vanishes if and only if $z = y$.
Note that the MSE loss function is a special case of the Bregman divergence with $\phi(z) = \| z \|_2^2$ and that the Poisson negative log-likelihood loss function is also a special case of the Bregman divergence with $\phi(z) = \sum_i z[i] \log z[i] - z[i]$.

Due to our conventions, the $z$ argument denotes the prediction and the $y$ argument denotes the target.
As the two variables are not symmetric in the Bregman divergence definition, we can swap the role played by $z$ and $y$ to get a different interpretation of the Bregman divergence.
In the target-first case, we have :
\begin{align*}
    \nabla_z \mathbb{E}[D_\phi(Y, z) | X=x] &= \nabla_z \mathbb{E}[\phi(Y) - \phi(z) - \langle \nabla \phi(z), Y - z \rangle | X=x] \\
    &= \nabla^2 \phi(z) \left( z - \mathbb{E}[Y | X=x] \right)
\end{align*}
Therefore, under strict convexity of $\phi$, the minimizer of $\mathbb{E}[D_\phi(Y, z) | X=x]$ is given by $z^* = \mathbb{E}[Y | X=x]$.
The target-first Bregman loss elicits the conditional mean of $Y$ given $X=x$. The same reasoning gives for the prediction-first Bregman loss :
\begin{align*}
    \nabla_z \mathbb{E}[D_\phi(z, Y) | X=x] &= \nabla_z \mathbb{E}[\phi(z) - \phi(Y) - \langle \nabla \phi(Y), z - Y \rangle | X=x] \\
    &= \nabla \phi(z) z - \mathbb{E}[\nabla \phi(Y) Y | X=x]
\end{align*}
which yields the minimizer $z^* = (\nabla \phi)^{-1}(\mathbb{E}[\nabla \phi(Y) | X=x])$, hence the quasi-arithmetic mean in the transformed space defined by $\nabla \phi$.
We can now state the following propositions regarding the Noise2Noise framework in the context of Bregman geometry.

\begin{proposition}
    \label{prop:b_target_first}
    If noise is centered, i.e. $\mathbb{E}[Y | X=x] = x$, then minimizing the Noise2Noise loss with a \textbf{target-first} Bregman divergence yields the same estimator as minimizing the supervised loss as $N \to \infty$.
\end{proposition}
\begin{proof}
    We derive the noise-induced loss bias in the context of Bregman geometry with target-first Bregman divergence :
    \begin{align*}
        \Delta_\rho(z, x) &= \mathbb{E}[D_\phi(Y, z) | X=x] - D_\phi(x, z) \\
        &= \mathbb{E}[\phi(Y) - \phi(z) - \langle \nabla \phi(z), Y - z \rangle | X=x] - (\phi(x) - \phi(z) - \langle \nabla \phi(z), x - z \rangle) \\
        &= \mathbb{E}[\phi(Y) | X=x] - \phi(x) + \langle \nabla \phi(z), x - \mathbb{E}[Y | X=x] \rangle
    \end{align*}
The term in $z$ vanishes as $\mathbb{E}[Y | X=x] = x$.
\end{proof}
\noindent Proposition \ref{prop:b_target_first} states that any target-first Bregman divergence suits the Noise2Noise framework, but the performance of the resulting estimator will depend on the choice of the Bregman divergence which has to be adapted to the noise model.
\begin{proposition}
    Minimizing the Noise2Noise loss with a \textbf{prediction-first} Bregman divergence yields the same estimator as minimizing the supervised loss as $N \to \infty$ if and only if :
    \begin{equation}
        \label{prop:b_prediction_first}
        \mathbb{E}[\nabla \phi(Y) | X=x] = \nabla \phi(x)
    \end{equation}
    If we also assume that the noise is centered, i.e, $\mathbb{E}[Y | X=x] = x$, then this condition is equivalent to $\phi$ being a quadratic function :
    \begin{equation}
        \phi(z) = \frac{1}{2} z^T Q z + b^T z + c
    \end{equation}
    with $Q$ a positive definite matrix, $b \in \mathbb{R}^d$ and $c \in \mathbb{R}$.
\end{proposition}

\begin{proof}
We derive the noise-induced loss bias in the context of Bregman geometry with prediction-first Bregman divergence :
\begin{align*}
    \Delta_\rho(z, x) &= \mathbb{E}[D_\phi(z, Y) | X=x] - D_\phi(z, x) \\
    &= \mathbb{E}[\phi(z) - \phi(Y) - \langle \nabla \phi(Y), z - Y \rangle | X=x] - (\phi(z) - \phi(x) - \langle \nabla \phi(x), z - x \rangle) \\
    &= -\mathbb{E}[\phi(Y) | X=x] + \phi(x) - \mathbb{E}[\langle \nabla \phi(Y), z - Y \rangle | X=x] + \langle \nabla \phi(x), z - x \rangle
\end{align*}
Let $g_x(z) = \mathbb{E}[\langle \nabla \phi(Y), z - Y \rangle | X=x] - \langle \nabla \phi(x), z - x \rangle$, we have $g_x'(z) = \mathbb{E}[\nabla \phi(Y) | X=x] - \nabla \phi(x)$.
In order to apply Corollary \ref{coro:delta}, we need $g_x'(z) = 0$ for all $z$, which is equivalent to $\mathbb{E}[\nabla \phi(Y) | X=x] = \nabla \phi(x)$.
It remains to be shown that this condition is equivalent to the gradient of $\phi$ being affine.
If noise is centered, then we can write $\mathbb{E}[\nabla \phi(Y) | X=x] = \nabla \phi(\mathbb{E}[Y | X=x])$.
We denote $g(z) = \nabla \phi(z)$, then the condition becomes $g(\mathbb{E}[Y | X=x]) = \mathbb{E}[g(Y) | X=x]$.
If $g$ is affine, the result is immediate. Conversely, if the condition holds for all $x$, then we can write $g(\mathbb{E}[Y | X=x]) = \mathbb{E}[g(Y) | X=x]$ for any random variable $Y$ with finite expectation.
For a two-point distribution $Y \sim p \delta_{y_1} + (1-p) \delta_{y_2}$, we have $\mathbb{E}[Y] = p y_1 + (1-p) y_2$ and $\mathbb{E}[g(Y)] = p g(y_1) + (1-p) g(y_2)$, hence $g(p y_1 + (1-p) y_2) = p g(y_1) + (1-p) g(y_2)$ for all $y_1, y_2$ and $p \in [0, 1]$, which is the definition of an affine function.
By integrating, we obtain that $\phi(x) = \frac{1}{2} x^T Q x + b^T x + c$ and the strict convexity of $\phi$ implies that $Q$ is positive definite.
\end{proof}

\subsubsection{Mahalanobis distance and correlated noise}

Proposition \ref{prop:b_prediction_first} can be interpreted as a generalization of the MSE loss function to the case of correlated noise.
Plugging $\phi(z) = \frac{1}{2} z^T Q z$ into the Bregman divergence definition \eqref{bregman_divergence} gives us :
\begin{align}
    D_\phi(z, y) &= \frac{1}{2} z^T Q z - \frac{1}{2} y^T Q y - \langle Q y, z - y \rangle \\
    &= \frac{1}{2} z^T Q z - \frac{1}{2} y^T Q y - y^T Q z + y^T Q y \\
    &= \frac{1}{2} (z - y)^T Q (z - y)
\end{align}
\noindent which is the Mahalanobis distance between $z$ and $y$ with respect to the positive definite matrix $Q$.
This is basically a generalization of Proposition \ref{prop:noise2noise_centered_mse}.
This means that any weighted MSE loss function with a positive definite weight matrix $Q$ is appropriate for the Noise2Noise framework under centered noise.
This has some potential benefits in the case of correlated noise where the covariance matrix of the noise is not diagonal, as the Mahalanobis distance can be used to decorrelate the noise and improve the denoising performance.
One can also choose diagonal weights to give more importance to some pixels than others, for example in the case of a region of interest in the image or sensitivity maps in medical imaging.




% \subsection{Other Noise2Noise-based frameworks}


% \subsubsection{Noise2Noise}

% In the context of a known noise model, the Noise2Noise principle can be extended to approaches like \textbf{Noisier2Noise} \cite{moran_noisier2noise_2019} where we can add extra noise to a single noisy image to create a pair of noisy images with a lower noise level in the target image.

% \subsubsection{Noise2Score}

% \cite{kim_noise2score_nodate} propose a 2-step framework called \textbf{Noise2Score}. First the score function is being estimated using a DAE trained under framework closed to Noisier2Noise where noise is embedded into the images according to the noise model and then the network learn the mapping to retrieve less noisy images.
% Then, the estimated denoised images are simply computed using the appropriate Tweedie's formula variant.
% This method is very promissing at it is at the crossroad between bayesian methods with strong mathematical background and self-supervised learning.


% \subsubsection{Noise2NoiseFlow}

% \cite{maleky_noise2noiseflow_2022} propose a framework at the crossroad between density estimation and the Noise2Noise framework. Density estimation is performed using Noise Flow \cite{abdelhamed_noise_2019}, a normalizing flow model designed to learn complex noise distributions from data.
% Notice that the density is closely related to the score.
% Whereas classical score matching methods require model trained in a supervised manner with clean data or denoising auto-encoders trained in an unsupervised manner to estimate the score function, they propose to directly learn the score function from noisy data only using the Noise2Noise principle.
% Both the denoising model and the normalizing flow are trained jointly using pairs of noisy images under the following loss function :

% \begin{equation}
%     \mathcal{L}(\theta, \phi) =  \underset{\substack{y_1, y_2 \sim p_{noisy}}}{\mathbb{E}} \left[ - \log p_{N_\phi}(y_1 | f_\theta(y_2)) - \log p_{N_\phi}(y_2 | f_\theta(y_1)) \right] + \lambda \mathcal{L}_{n2n}(\theta)
% \end{equation}

% with $p_{N_\phi}$ being the noise distribution modeled by the normalizing flow with parameters $\phi$, $f_\theta$ the denoising model with parameters $\theta$ and $\mathcal{L}_{n2n}(\theta)$ the standard Noise2Noise loss. The authors of Noise2NoiseFlow show that this additional term gives more stability to the training process and leads to better denoising results.

% Note that this framework is useless if we already know the natural noise distribution in the data as we could directly use it to denoise images using bayesian methods.

\section{Bootstrapping noisy measurements}

In this section, we discuss how to generate two independant measurements to apply Noise2Noise on from a single measurement.

\subsection{Statistical bootstrapping}

In the context of additive gaussian noise, \textbf{Recorrupted2Recorrupted} \cite{Pang2021} aims to provide a way to simulate noisy image pairs satisfying statistical conditions from single measurements.
This can be particularly usefull in the case of medical imaging where we often have only one noisy image per patient and where the noise model is well-known.
We propose an intuitive alternative to their method with Poisson noise.

\begin{definition}{Infinitely divisible distribution}
    Let $Y$ be a random variable. $Y$ is said to be $n$-divisible if there exists a sequence of $n$ i.i.d. random variables $\{Y_i\}_{i=1}^n$ such that $Y \stackrel{d}{=} \sum_{i=1}^n Y_i$. If $Y$ is $n$-divisible for all $n \in \mathbb{N}$, then $Y$ is said to be infinitely divisible.
\end{definition}

If the noise model is well-known and satisfies the infinite divisibility property, we can use this property to generate pairs of noisy images from a single noisy image by splitting the noise into two independent parts.

\begin{proposition}
    The Poisson distribution is infinitely divisible ; in particular, it is $2$-divisible. Thus, if $Y \sim \text{Poisson}(\lambda)$, then there exists two i.i.d. random variables $Y_1$ and $Y_2$ such that $Y = Y_1 + Y_2$ and $Y_1, Y_2 \sim \text{Poisson}(\lambda/2)$.
    \label{prop:poisson_divisibility}
\end{proposition}

\begin{proof}
    This proof will focus on the case of $2$-divisibility, but it can be easily extended to the general case of $n$-divisibility.
    Let $Y$ be a random variable following a Poisson distribution with parameter $\lambda$, and let $Y_1 | Y=y \sim \text{Binomial}(y, 1/2)$ be a random variable following a binomial distribution with parameters $y$ and $1/2$ conditionally on $Y=y$. We define $Y_2 = Y - Y_1$. We have the following result for the distribution of $Y_1$ :
    \begin{align*}
        \mathbb{P}(Y_1 = k) &= \sum_{y=k}^{\infty} \mathbb{P}(Y_1 = k | Y = y) \mathbb{P}(Y = y) \\
        &= \sum_{y=k}^{\infty} \binom{y}{k} \left( \frac{1}{2} \right)^y \frac{\lambda^y e^{-\lambda}}{y!} \\
        &= \sum_{y=k}^{\infty} \frac{y!}{k!(y-k)!} \left( \frac{1}{2} \right)^y \frac{\lambda^y e^{-\lambda}}{y!} \\
        &= \frac{1}{k!} \left( \frac{\lambda}{2} \right)^k e^{-\lambda/2} \sum_{y=k}^{\infty} \frac{(\lambda/2)^{y-k} e^{-\lambda/2}}{(y-k)!} \\
        &= \frac{1}{k!} \left( \frac{\lambda}{2} \right)^k e^{-\lambda/2}
    \end{align*}
    Thus, $Y_1$ follows a Poisson distribution with parameter $\lambda/2$. By symmetry, $Y_2$ also follows a Poisson distribution with parameter $\lambda/2$. Moreover, we have $Y = Y_1 + Y_2$ by definition.

    Finally, we show that $Y_1$ and $Y_2$ are independent. We have :
    \begin{align*}
        \mathbb{P}(Y_1 = k, Y_2 = l) &= \mathbb{P}(Y_1 = k, Y_2 = l | Y = k + l) \mathbb{P}(Y = k + l) \\
        &= \binom{k+l}{k} \left( \frac{1}{2} \right)^{k+l} \frac{\lambda^{k+l} e^{-\lambda}}{(k+l)!} \\
        &= \frac{1}{k!} \left( \frac{\lambda}{2} \right)^k e^{-\lambda/2} \cdot \frac{1}{l!} \left( \frac{\lambda}{2} \right)^l e^{-\lambda/2} \\
        &= \mathbb{P}(Y_1 = k) \mathbb{P}(Y_2 = l)
    \end{align*}
\end{proof}

The previous proposition shows that it is possible to perform binomial thinnning on Poisson data in order to generate independant pairs of noisy images.
However, subsampled images expectation is not the same as the original image expectation, as well as noise variance, which means that the signal to noise ratio is higher in the subsampled images than in the original image.

\cite{gr2r} go beyond the additive Gaussian noise model or Poisson noise model and propose a generic formulation for any posterior distribution $p(Y|X)$ that belongs to the exponential family of distributions.
Their implementation in the case of Poisson noise is shown in the following proposition.
We introduce $\gamma$ and $\alpha$ which respectively control the noise model scale factor and the balance between the two measurements $Y_1$ and $Y_2$ :

\begin{equation}
    \begin{cases}
        Y = \gamma Z \\
        Z | X =x \sim \text{Poisson}(\frac{x}{\gamma}) \\
        W | Z = z \sim \text{Binomial}(z, \alpha) \\
        Y_1 = \frac{Y - \gamma W}{1 - \alpha} \\
        Y_2 = \frac{1}{\alpha} Y - \frac{1 - \alpha}{\alpha} Y_1
    \end{cases}
    \label{eq:gr2r}
\end{equation}

\begin{proposition}
    The measurements admits decomposition $Y = (1 - \alpha) Y_1 + \alpha Y_2$ and satisfy the following properties :
    \begin{itemize}
        \item $Y_1$ and $Y_2$ are independent given the clean image $X$.
        \item $\mathbb{E}[Y_1 | X=x] = \mathbb{E}[Y_2 | X=x] = \mathbb{E}[Y | X=x] = x$.
        \item $\mathbb{V}[Y_1 | X=x] = \frac{1}{1 - \alpha} \mathbb{V}[Y | X=x] = \frac{1}{1 - \alpha} x$.
        \item $\mathbb{V}[Y_2 | X=x] = \frac{1}{\alpha} \mathbb{V}[Y | X=x] = \frac{1}{\alpha} x$.
    \end{itemize}
\end{proposition}

\cite{gr2r} propose a generic proof for the previous proposition when the noise distribution is any member of the exponential family of distributions, which includes the Poisson distribution as a special case.

\begin{remark}
    The definition of $Y_1$ and $Y_2$ in \eqref{eq:gr2r} when $\alpha = 0.5$ and $\gamma = 2$ match the one from proposition \ref{prop:poisson_divisibility}.
\end{remark}

\subsection{Data-based bootstrapping}


\textbf{Noise2Void} \cite{krull_noise2void_2019} and \textbf{Noise2Self} \cite{batson_noise2self_nodate} frameworks further extend the idea of training denoising models without clean targets by using only single noisy images.
\cite{batson_noise2self_nodate} introduce the concept of \textit{$\mathcal{J}$-invariant} functions, which are functions that do not depend on certain subsets of their input.
We train a denoising model to be $\mathcal{J}$-invariant by masking certain pixels of the input image and predicting their values based on the surrounding pixels.
The work of \cite{krull_noise2void_2019} implements this idea by randomly masking pixels in the input image during training and using the unmasked pixels to predict the masked ones.
This is called the \textit{blind-spot} strategy and ensures that the model does not learn to simply copy the noisy input to the output, thus learning the identity.

In some applications, it is also possible to take advantage of the acquisition process to generate multiple noisy measurements of the same underlying clean image.
If a signal $y$ is drawn from \eqref{eq:inverse-problem}, it can be splitted into two non-overlapping subsets of measurements $y_1$ and $y_2$ such that

\begin{equation}
    \begin{cases}
        y_1 = A_1 x + b_1 \\
        y_2 = A_2 x + b_2
    \end{cases}
\end{equation}
with $A_1$ and $A_2$ two linear operators inheriting from $A$ and defined accordingly to $y_1$ and $y_2$ and $b_1$ and $b_2$ two noise vectors.

This idea was exploited in the context of Computed Tomography where odd and even projections are separated to get $K$ sinogram subsets \cite{hendriksen_noise2inverse_2020}.
Each resulting sinogram subset is an even more discretized version of the original sinogram, thus even more ill-posed and noisy than the original problem.
They study two different training strategies.
The first one uses the average of the first $K-1$ sinograms as input and the remaining one as target. This strategy is close to the Noise2Self principle.
The second one is the reversed of the first and is closer to the Noisier2Noise strategy as the input will be noisier than the target.


\section{Generalization of Noise2Noise to linear inverse problems}

In this section, we get back to the general framework where $A$ is unspecified.
Let again $\{(Y_i^{(1)}, Y_i^{(2)})\}_{i=1}^N$ be a finite sequence of independant and statistically distributed random variables and $y_i^{(j)} = A_i^{(j)}x_i + b_i^{(j)}$ the observed degraded images.
We also rewrite \ref{eq:inverse-problem} as :
\begin{equation}
    \begin{cases}
        Y^{(j)} | X=x \sim \mathcal{C}(A^{(j)}x), \quad j = 1, 2 \\
        Y^{(j)} = A^{(j)} X + B^{(j)}, \quad j = 1, 2
    \end{cases}
\end{equation}
with $A^{(j)} \in \mathbb{R}^{d \times m}$ two linear operators inheriting from $A$ and defined accordingly to $Y^{(1)}$ and $Y^{(2)}$.


In the context of non-blind estimation, i.e. when $A^{(j)}$ is known, we can optimize a neural network $f_{\theta}$ under the following objective function \cite{xia_training_nodate} :
\begin{equation}
    \mathcal{L}_{swap}(\theta) = \frac{1}{N} \sum_{i=1}^N \left[ \rho(A_i^{(2)} f_{\theta}(y_i^{(1)}) - y_i^{(2)}) + \rho(A_i^{(1)} f_{\theta}(y_i^{(2)}) - y_i^{(1)}) \right]
    \label{eq:xia_swap_loss}
\end{equation}


\begin{theorem}
    Minimizing \ref{eq:n2n_loss} with $\rho(x, y) = \| x - y \|_2^2$ gives the same estimator as minimizing the homologous supervised loss as $N \to \infty$ under Noise2Noise statistical conditions.
\end{theorem}

\subsubsection{MSE loss}

\begin{proof}
    Proof in theorem \ref{th:noise2noise_l2} is still valid as long as we replace $f_{\theta}(Y^{(j)})$ by $Af_{\theta}(Y^{(j)})$. We thus have :
    \begin{align*}
        R_{N}^{n2n}(\theta) &= \frac{1}{N} \sum_{i=1}^N \left[ \| A_i^{(2)} f_{\theta}(Y_i^{(1)}) - Y_i^{(2)} \|_2^2 + \| A_i^{(1)} f_{\theta}(Y_i^{(2)}) - Y_i^{(1)} \|_2^2 \right] \\
        &= 2 R_{N}^{clean}(\theta) + \frac{1}{N} \sum_{i=1}^N \left[ \| B_i^{(1)} \|_2^2 + \| B_i^{(2)} \|_2^2 \right] + o(1)
    \end{align*}
\end{proof}

\subsubsection{Poisson linear inverse problems}

With the conventions of section \ref{sec:noise2noise_poisson}, we now consider the case of linear inverse problems corrupted by Poisson noise :

\begin{equation}
    Y^{(j)} \sim Poisson(A^{(j)} X), \quad j = 1, 2
\end{equation}

\begin{theorem}
    If $Y | X=x \sim \text{Poisson}(Ax)$, then minimizing \ref{eq:n2n_loss} with $\rho = \mathcal{L}_{Poisson}$ gives the same estimator as minimizing the homologous supervised loss as $N \to \infty$ under Noise2Noise statistical conditions.
\end{theorem}

\begin{proof}
    Proof in theorem \ref{th:noise2noise_poisson} is still valid as long as we replace $f_{\theta}(Y^{(j)})$ by $A f_{\theta}(Y^{(j)})$. If we define the empirical risk as :
    \begin{equation*}
        R_{N}^{n2n}(\theta) = \frac{1}{N} \sum_{i=1}^N \left[ \mathcal{L}_{Poisson}(A_i^{(2)} f_{\theta}(Y_i^{(1)}), Y_i^{(2)}) + \mathcal{L}_{Poisson}(A_i^{(1)} f_{\theta}(Y_i^{(2)}), Y_i^{(1)}) \right]
    \end{equation*}
    then we have :
    \begin{equation*}
        R_{N}^{n2n}(\theta) \xrightarrow[N \to \infty]{a.s.} 2 \mathbb{E} \left[ f_\theta(Y) - A X \log(f_\theta(Y)) + \log(Y!)\right]
    \end{equation*}
\end{proof}

\subsection{Splitting strategies}

\subsubsection{Data splitting}

eg tachella, Noise2Inverse

\subsubsection{Statistical splitting}

eg R2R, GR2R

\subsection{Single-operator framework}

There exists a special case of inverse problems where $A^{(1)} = A^{(2)}$ for each training pair.
In this case, the loss function can be simplified as follows :

\begin{equation}
    \mathcal{L}_{\text{swap}}(\theta) = \frac{1}{N} \sum_{i=1}^N \left[ \rho(A_i f_{\theta}(y_i^{(1)}) - y_i^{(2)}) + \rho(A_i f_{\theta}(y_i^{(2)}) - y_i^{(1)}) \right]
\end{equation}

\subsection{Data consistency enforcement}

Eventhough the loss function \ref{eq:xia_swap_loss} is sufficient to guarantee that the optimal solution of the problem is the same as the one of the supervised problem, it does not enforce data-consistency during training, which can lead to suboptimal results.
Indeed, this term is biased toward smoothness as a denoising objective. Thus it can be usefull to add an extra term to the loss function to enforce data-consistency during training \cite{xia_training_nodate} :

\begin{equation}
    \mathcal{L}_{\text{self}} = \frac{1}{N} \sum_{i=1}^N \left[ \rho(A_i^{(1)} f_{\theta}(y_i^{(1)}) - y_i^{(1)}) + \rho(A_i^{(2)} f_{\theta}(y_i^{(2)}) - y_i^{(2)}) \right]
\end{equation}

In the case of inifinite divisibility-based frameworks where $y = y_1 + y_2$ and $y_1$ and $y_2$ are derived from the same probability law, we can simply use the following data-consistency term :

\begin{equation}
    \mathcal{L}_{\text{self}} = \frac{1}{N} \sum_{i=1}^N \rho\left( \frac{A_i f_{\theta}(y_i^{(1)}) + A_i f_{\theta}(y_i^{(2)})}{2} - y\right)
\end{equation}
where we want the resulting average to match the original noisy image $y$.

Eitherway, the final loss function is then written as :
\begin{equation}
    \mathcal{L}_{\text{total}} = \mathcal{L}_{\text{swap}} + \lambda \mathcal{L}_{\text{self}}
\end{equation}
with $\lambda$ a hyperparameter to tune the weight of the data-consistency term.